% Generated draft: requires mathematical and source review.
\documentclass[11pt]{article}
\usepackage[margin=25mm]{geometry}
\usepackage{fontspec}
\setmainfont{Latin Modern Roman}
\usepackage{amsmath,amssymb,mathtools,unicode-math}
\setmathfont{Latin Modern Math}
\usepackage{enumitem,needspace,xurl,hyperref}
\setlength{\parindent}{0pt}
\setlength{\parskip}{5pt}
\setlength{\emergencystretch}{3em}
\newcommand{\N}{\mathbb{N}}
\newcommand{\Z}{\mathbb{Z}}
\newcommand{\Q}{\mathbb{Q}}
\newcommand{\R}{\mathbb{R}}
\newcommand{\C}{\mathbb{C}}
\newcommand{\F}{\mathbb{F}}
\newcommand{\A}{\mathbb{A}}
\newcommand{\PP}{\mathbb{P}}
\newcommand{\E}{\mathbb{E}}
\newcommand{\eps}{\varepsilon}
\newcommand{\dd}{\,\mathrm{d}}
\begin{document}
\section*{UnsolvedMath ID: 30006658}
\textbf{Source problem code:} catalogue entry
\subsection{Definitions and mathematical statement}
Let \(p\) be a prime and let \(G\) be a compact Lie group of dimension \(d>0\). Write \(G^\circ\) for its identity component and \(\pi_0(G)=G/G^\circ\) for its finite group of components. A finite group is \(p\)-nilpotent if it has a normal subgroup of order prime to \(p\) and index a power of \(p\). Assume that \(\pi_0(G)\) is not \(p\)-nilpotent. Write \(BG\) for the classifying space of \(G\), \(BG^\wedge_p\) for its \(p\)-completion, and \(\Omega(BG^\wedge_p)\) for the based loop space of that completion.

Retain the source notation
\[
A=C_*(\Omega(BG^\wedge_p)),\qquad
\varepsilon\colon\Gamma_k A\longrightarrow A
\]
for the chain object and the map underlying the displayed question. Here \(\varepsilon\) is only a label for that map. The supplied record does not specify the coefficient object \(k\), the coefficient convention for \(C_*\), the definition of \(\Gamma_k\), or the construction of this map. These are unresolved parts of the formulation; in particular, no particular field or cellularization convention is assumed here.

\textbf{Question (conditional on those conventions being supplied).} For every such \(G\), is the induced graded map
\[
\pi_*(\varepsilon)\colon
\pi_*(\Gamma_k C_*(\Omega(BG^\wedge_p)))
\longrightarrow
\pi_*(C_*(\Omega(BG^\wedge_p)))
\]
zero? Equivalently, in the intended spectral interpretation, is
\[
\pi_n(\varepsilon)=0\qquad\text{for every }n\in\Z?
\]
The intended range of coefficient objects \(k\) and the precise homotopical setting remain unspecified, so this is a conditional transcription rather than a fully determined mathematical formulation.

The accompanying excerpt says that a suspension makes a map denoted \(\nu_*\) potentially nonzero in degrees from \(0\) through \(d\), and reports that Benson's squeezed resolutions show that, for finite groups, \(\nu_*\) is an isomorphism in the \(p\)-nilpotent case and zero otherwise. These are supplied contextual claims. The excerpt does not define \(\nu_*\) or explicitly identify it with the displayed map, so neither the identification nor a restriction of the question to degrees \(0,\ldots,d\) is imposed here.
\subsection{Short English statement}
The question asks whether the displayed map induces zero on all homotopy groups whenever \(G\) is a positive-dimensional compact Lie group whose component group is not \(p\)-nilpotent. The supplied record labels this question open. Its coefficient object \(k\), functor \(\Gamma_k\), and underlying map require clarification before the statement is fully precise.
\subsection{Sources}
\begin{itemize}
\item Local UnsolvedMath registry, \texttt{data/unsolvedmath/problems.json}, permanent ID \(30006658\), source code \texttt{catalogue entry}; supplied complete source record, titled ``Vanishing for Non-\(p\)-Nilpotent Compact Lie Groups.'' The original excerpt labels the question ``Question 2.2.''
\item \emph{Cohomology of Finite Groups: Interactions and Applications} (2026), Oberwolfach Reports, pp.\ 541--610, DOI \texttt{10.4171/owr/2026/9}. Supplied source link: \url{https://doi.org/10.4171/owr/2026/9}.
\item Supplied literature reference: Benson, ``The Omega-Tate spectrum,'' Oberwolfach Report 9/2026, \url{https://doi.org/10.4171/owr/2026/9}. The record does not separately identify a proposer.
\item Literature assessment dated exactly \texttt{2026-08-22}: the supplied curation labels the compact-Lie-group question open, reports the finite-group analogue as known, and says that no proof or counterexample for the stated general case was found. This is supplied dated curation, not an independently verified literature assessment.
\item The registry labels the statement \texttt{corrected\_verified} and states that it was reconstructed and checked against the official report. That is supplied provenance, not a verification performed for this draft. The excerpt also cites Benson's squeezed resolutions as reference [2], without supplying the corresponding bibliographic entry.
\end{itemize}
\end{document}
